\documentclass[a4paper,twoside]{article}

\usepackage{morewrites}

\usepackage[svgnames,dvipsnames,x11names]{xcolor}
\usepackage{graphicx}
\usepackage{texsmith-lists}
\usepackage[english]{babel}
\usepackage[colorlinks=true, linkcolor=NavyBlue, urlcolor=NavyBlue, citecolor=NavyBlue]{hyperref}
\usepackage[a4paper]{geometry}
\usepackage{ts-fonts}
\usepackage{tabularx}
\usepackage{amsmath}
\usepackage{float}
\usepackage{seqsplit}
\newcolumntype{Y}{>{\centering\arraybackslash}X}
\newcolumntype{Z}{>{\raggedleft\arraybackslash}X}
\usepackage{ts-code}

\usepackage{lastpage}
\usepackage{refcount}
\makeatletter
\def\ps@plain{
\let\@oddhead\@empty
\let\@evenhead\@empty
\def\@oddfoot{
\hfil
\ifnum\getpagerefnumber{LastPage}>1\relax
\thepage
\fi
\hfil}
\let\@evenfoot\@oddfoot
}
\makeatother

\title{Tables}
\author{}\date{}
\begin{document}
\maketitle
\textbf{Mostly settled.} The answer TeXSmith shipped is the \tscodeinline{yaml table} fence,
which the TMark parser reads and validates: grouped headers, column widths
(fixed, auto, relative), a \tscodeinline{width-\allowbreak{}group} that makes several columns share one
width, per-column and per-group alignment, row and column spans (the absorbed
cells written as \tscodeinline{\textasciitilde{}}), separators, and a caption line after the fence. The
table model is specified and round-trips through \tscodeinline{tmark fmt}, and a shape it
cannot express is rejected with a \tscodeinline{table-\allowbreak{}*} diagnostic rather than silently
flattened. See Tables.

What is still open from the list below is table \textbf{orientation} (rotating a
very large table) and the width-control knobs (\tscodeinline{tabulary} and friends); both
are roadmap items.
\section{Complex Tables}
Markdown offers limited table configuration---only column alignment by default. PyMdown provides captions, and superfences can inject more metadata, but we still miss:
\begin{itemize}
\item Table width (auto vs. full width)
\item Resizing oversized tables
\item Orientation (portrait vs. landscape)
\item Column and row spanning
\item Horizontal and vertical separator lines
\item Column widths (fixed, auto, relative)
\end{itemize}
\subsection{Extended Markdown Table Syntax}
An idea that was considered and not taken: extending the pipe-table syntax
itself so a cell could span, with \tscodeinline{>>>} spanning horizontally and \tscodeinline{vvv}
spanning vertically.
\begin{tscode}[lang=markdown, engine=pygments]
\begin{Verbatim}[commandchars=\\\{\},breaklines, breakanywhere, commandchars=\\\{\}]
| Header 1 | Header 2 | Header 3 |
|:\PYZhy{}\PYZhy{}\PYZhy{}\PYZhy{}\PYZhy{}\PYZhy{}\PYZhy{}\PYZhy{}\PYZhy{}|:\PYZhy{}\PYZhy{}\PYZhy{}\PYZhy{}\PYZhy{}\PYZhy{}\PYZhy{}\PYZhy{}:|\PYZhy{}\PYZhy{}\PYZhy{}\PYZhy{}\PYZhy{}\PYZhy{}\PYZhy{}\PYZhy{}\PYZhy{}:|
| Cell 1   | Cell 2   | Cell 3   |
| Cell 4   | \PYZgt{}\PYZgt{}\PYZgt{}      | Cell 6   |
| Cell 7   | Cell 8   | Cell 11  |
| Cell 9   |          | vvv      |
\end{Verbatim}
\end{tscode}
\subsection{Cmidrule example}
\begin{tscode}[lang=latex, engine=pygments]
\begin{Verbatim}[commandchars=\\\{\},breaklines, breakanywhere, commandchars=\\\{\}]
\PY{k}{\PYZbs{}begin}\PY{n+nb}{\PYZob{}}tabular\PY{n+nb}{\PYZcb{}}\PY{n+nb}{\PYZob{}}@\PY{n+nb}{\PYZob{}}\PY{n+nb}{\PYZcb{}}lll@\PY{n+nb}{\PYZob{}}\PY{n+nb}{\PYZcb{}}\PY{n+nb}{\PYZcb{}}
\PY{k}{\PYZbs{}toprule}
\PY{n+nb}{\PYZam{}} \PY{k}{\PYZbs{}multicolumn}\PY{n+nb}{\PYZob{}}2\PY{n+nb}{\PYZcb{}}\PY{n+nb}{\PYZob{}}c\PY{n+nb}{\PYZcb{}}\PY{n+nb}{\PYZob{}}Reaction\PY{n+nb}{\PYZcb{}} \PY{k}{\PYZbs{}\PYZbs{}}
\PY{k}{\PYZbs{}cmidrule}(l)\PY{n+nb}{\PYZob{}}2\PYZhy{}3\PY{n+nb}{\PYZcb{}}
Person \PY{n+nb}{\PYZam{}} Face \PY{n+nb}{\PYZam{}} Exclamation \PY{k}{\PYZbs{}\PYZbs{}}
\PY{k}{\PYZbs{}midrule}
\PY{k}{\PYZbs{}multirow}\PY{n+na}{[t]}\PY{n+nb}{\PYZob{}}2\PY{n+nb}{\PYZcb{}}\PY{n+nb}{\PYZob{}}*\PY{n+nb}{\PYZcb{}}\PY{n+nb}{\PYZob{}}VIPs\PY{n+nb}{\PYZcb{}} \PY{n+nb}{\PYZam{}} :) \PY{n+nb}{\PYZam{}} Nice \PY{k}{\PYZbs{}\PYZbs{}}
\PY{n+nb}{\PYZam{}} :] \PY{n+nb}{\PYZam{}} Sleek \PY{k}{\PYZbs{}\PYZbs{}}
Others \PY{n+nb}{\PYZam{}} \PY{k}{\PYZbs{}multicolumn}\PY{n+nb}{\PYZob{}}2\PY{n+nb}{\PYZcb{}}\PY{n+nb}{\PYZob{}}c\PY{n+nb}{\PYZcb{}}\PY{n+nb}{\PYZob{}}Not available\PY{n+nb}{\PYZcb{}} \PY{k}{\PYZbs{}\PYZbs{}}
\PY{k}{\PYZbs{}bottomrule}
\PY{k}{\PYZbs{}end}\PY{n+nb}{\PYZob{}}tabular\PY{n+nb}{\PYZcb{}}
\end{Verbatim}
\end{tscode}
\subsection{Align to decimal points}
Find a syntax to align numbers to dot. \tscodeinline{lS@\{\}}
\subsection{Raw table syntax}
Superfences do not work directly with tables, so define a \tscodeinline{table} fence that accepts YAML options:
\begin{tscode}[lang=markdown, engine=pygments]
\begin{Verbatim}[commandchars=\\\{\},breaklines, breakanywhere, commandchars=\\\{\}]
\PY{l+s+sb}{```}\PY{l+s+sb}{table}
\PY{l+s}{width: full}
\PY{l+s}{caption: Sample Table}
\PY{l+s}{label: tbl:sample}
\PY{l+s}{header: true}
\PY{l+s}{print:}
\PY{l+s}{  orientation: landscape}
\PY{l+s}{  resize: true}
\PY{l+s}{columns:}
\PY{l+s}{  \PYZhy{} alignment: left}
\PY{l+s}{    width: auto}
\PY{l+s}{  \PYZhy{} alignment: center}
\PY{l+s}{    width: 2cm}
\PY{l+s}{  \PYZhy{} alignment: right}
\PY{l+s}{    width: auto}
\PY{l+s}{rows:}
\PY{l+s}{  \PYZhy{} height: auto}
\PY{l+s}{    alignment: top}
\PY{l+s}{    span: [2, auto]}
\PY{l+s}{data:}
\PY{l+s}{  \PYZhy{} [Cell 1, Cell 2, Cell 3]}
\PY{l+s}{  \PYZhy{} [Cell 4, null, Cell 6]}
\PY{l+s}{  \PYZhy{} [Cell 7, Cell 8, null]}
\PY{l+s+sb}{```}
\end{Verbatim}
\end{tscode}
The goal is to convert Markdown tables into \LaTeX{} tables with automatic line breaks so that columns wrap gracefully when a table is too wide.
\section{Shrink-to-fit vs Auto-expand}
Long tables et auto-ajustement.
\begin{tscode}[lang=latex, engine=pygments]
\begin{Verbatim}[commandchars=\\\{\},breaklines, breakanywhere, commandchars=\\\{\}]
\PY{k}{\PYZbs{}documentclass}\PY{n+nb}{\PYZob{}}article\PY{n+nb}{\PYZcb{}}
\PY{k}{\PYZbs{}usepackage}\PY{n+na}{[margin=2cm]}\PY{n+nb}{\PYZob{}}geometry\PY{n+nb}{\PYZcb{}}
\PY{k}{\PYZbs{}usepackage}\PY{n+nb}{\PYZob{}}booktabs\PY{n+nb}{\PYZcb{}}
\PY{k}{\PYZbs{}usepackage}\PY{n+na}{[french]}\PY{n+nb}{\PYZob{}}babel\PY{n+nb}{\PYZcb{}}
\PY{k}{\PYZbs{}usepackage}\PY{n+nb}{\PYZob{}}array\PY{n+nb}{\PYZcb{}}

\PY{c}{\PYZpc{} ltablex : Le pont entre tabularx et longtable}
\PY{k}{\PYZbs{}usepackage}\PY{n+nb}{\PYZob{}}ltablex\PY{n+nb}{\PYZcb{}}

\PY{c}{\PYZpc{} IMPORTANT : On NE met PAS \PYZbs{}keepXColumns ici.}
\PY{c}{\PYZpc{} Sans cette commande, ltablex va calculer si le tableau a besoin}
\PY{c}{\PYZpc{} de toute la largeur ou non.}

\PY{k}{\PYZbs{}begin}\PY{n+nb}{\PYZob{}}document\PY{n+nb}{\PYZcb{}}

\PY{k}{\PYZbs{}section}\PY{k}{*}\PY{n+nb}{\PYZob{}}Cas 1 : Table petite (Compacte)\PY{n+nb}{\PYZcb{}}
\PY{c}{\PYZpc{} Ici, comme le texte est court, le tableau ne prendra pas toute la page}
\PY{c}{\PYZpc{} Les colonnes X vont se comporter comme des colonnes \PYZsq{}l\PYZsq{}}
\PY{k}{\PYZbs{}begin}\PY{n+nb}{\PYZob{}}tabularx\PY{n+nb}{\PYZcb{}}\PY{n+nb}{\PYZob{}}\PY{k}{\PYZbs{}textwidth}\PY{n+nb}{\PYZcb{}}\PY{n+nb}{\PYZob{}}lXX\PY{n+nb}{\PYZcb{}}
    \PY{k}{\PYZbs{}toprule}
    \PY{k}{\PYZbs{}textbf}\PY{n+nb}{\PYZob{}}ID\PY{n+nb}{\PYZcb{}} \PY{n+nb}{\PYZam{}} \PY{k}{\PYZbs{}textbf}\PY{n+nb}{\PYZob{}}Statut\PY{n+nb}{\PYZcb{}} \PY{n+nb}{\PYZam{}} \PY{k}{\PYZbs{}textbf}\PY{n+nb}{\PYZob{}}Note\PY{n+nb}{\PYZcb{}} \PY{k}{\PYZbs{}\PYZbs{}}
    \PY{k}{\PYZbs{}midrule}
    \PY{k}{\PYZbs{}endfirsthead}
    1 \PY{n+nb}{\PYZam{}} OK \PY{n+nb}{\PYZam{}} R.A.S. \PY{k}{\PYZbs{}\PYZbs{}}
    \PY{k}{\PYZbs{}midrule}
\PY{k}{\PYZbs{}end}\PY{n+nb}{\PYZob{}}tabularx\PY{n+nb}{\PYZcb{}}

\PY{k}{\PYZbs{}vspace}\PY{n+nb}{\PYZob{}}2cm\PY{n+nb}{\PYZcb{}}

\PY{k}{\PYZbs{}section}\PY{k}{*}\PY{n+nb}{\PYZob{}}Cas 2 : Table large (Extension automatique)\PY{n+nb}{\PYZcb{}}
\PY{c}{\PYZpc{} Ici, le texte est long. Le tableau va détecter qu\PYZsq{}il dépasse,}
\PY{c}{\PYZpc{} s\PYZsq{}étendre jusqu\PYZsq{}à \PYZbs{}textwidth, et activer le retour à la ligne.}
\PY{k}{\PYZbs{}begin}\PY{n+nb}{\PYZob{}}tabularx\PY{n+nb}{\PYZcb{}}\PY{n+nb}{\PYZob{}}\PY{k}{\PYZbs{}textwidth}\PY{n+nb}{\PYZcb{}}\PY{n+nb}{\PYZob{}}lXX\PY{n+nb}{\PYZcb{}}
    \PY{k}{\PYZbs{}toprule}
    \PY{k}{\PYZbs{}textbf}\PY{n+nb}{\PYZob{}}ID\PY{n+nb}{\PYZcb{}} \PY{n+nb}{\PYZam{}} \PY{k}{\PYZbs{}textbf}\PY{n+nb}{\PYZob{}}Description\PY{n+nb}{\PYZcb{}} \PY{n+nb}{\PYZam{}} \PY{k}{\PYZbs{}textbf}\PY{n+nb}{\PYZob{}}Analyse\PY{n+nb}{\PYZcb{}} \PY{k}{\PYZbs{}\PYZbs{}}
    \PY{k}{\PYZbs{}midrule}
    \PY{k}{\PYZbs{}endfirsthead}

    \PY{k}{\PYZbs{}textbf}\PY{n+nb}{\PYZob{}}ID\PY{n+nb}{\PYZcb{}} \PY{n+nb}{\PYZam{}} \PY{k}{\PYZbs{}textbf}\PY{n+nb}{\PYZob{}}Description\PY{n+nb}{\PYZcb{}} \PY{n+nb}{\PYZam{}} \PY{k}{\PYZbs{}textbf}\PY{n+nb}{\PYZob{}}Analyse\PY{n+nb}{\PYZcb{}} \PY{k}{\PYZbs{}\PYZbs{}}
    \PY{k}{\PYZbs{}midrule}
    \PY{k}{\PYZbs{}endhead}

    204 \PY{n+nb}{\PYZam{}}
    Ici j\PYZsq{}ai un texte suffisamment long pour justifier que le tableau prenne toute la largeur disponible sur la page. \PY{n+nb}{\PYZam{}}
    Et ici une autre colonne qui va se partager l\PYZsq{}espace restant équitablement avec la colonne précédente. \PY{k}{\PYZbs{}\PYZbs{}}

    205 \PY{n+nb}{\PYZam{}} Test de remplissage \PY{n+nb}{\PYZam{}} Encore du texte... \PY{k}{\PYZbs{}\PYZbs{}}
    \PY{k}{\PYZbs{}bottomrule}
\PY{k}{\PYZbs{}end}\PY{n+nb}{\PYZob{}}tabularx\PY{n+nb}{\PYZcb{}}
\PY{k}{\PYZbs{}end}\PY{n+nb}{\PYZob{}}document\PY{n+nb}{\PYZcb{}}
\end{Verbatim}
\end{tscode}
We could use a tester to have a rough idea if the table will fit or not, and decide to use \tscodeinline{tabularx} or \tscodeinline{ltablex} accordingly.
\begin{tscode}[lang=tex, engine=pygments]
\begin{Verbatim}[commandchars=\\\{\},breaklines, breakanywhere, commandchars=\\\{\}]
\PY{k}{\PYZbs{}documentclass}\PY{n+nb}{\PYZob{}}article\PY{n+nb}{\PYZcb{}}
\PY{k}{\PYZbs{}usepackage}\PY{n+nb}{\PYZob{}}booktabs\PY{n+nb}{\PYZcb{}}
\PY{k}{\PYZbs{}usepackage}\PY{n+nb}{\PYZob{}}tabularx\PY{n+nb}{\PYZcb{}} \PY{c}{\PYZpc{} ou tabularray (tblr), etc.}

\PY{k}{\PYZbs{}newsavebox}\PY{n+nb}{\PYZob{}}\PY{k}{\PYZbs{}tblbox}\PY{n+nb}{\PYZcb{}}

\PY{k}{\PYZbs{}begin}\PY{n+nb}{\PYZob{}}document\PY{n+nb}{\PYZcb{}}
\PY{k}{\PYZbs{}typeout}\PY{n+nb}{\PYZob{}}TEXTWIDTH\PY{n+nb}{\PYZus{}}PT=\PY{k}{\PYZbs{}the}\PY{k}{\PYZbs{}textwidth}\PY{n+nb}{\PYZcb{}}
\PY{k}{\PYZbs{}typeout}\PY{n+nb}{\PYZob{}}LINEWIDTH\PY{n+nb}{\PYZus{}}PT=\PY{k}{\PYZbs{}the}\PY{k}{\PYZbs{}linewidth}\PY{n+nb}{\PYZcb{}}
\PY{k}{\PYZbs{}typeout}\PY{n+nb}{\PYZob{}}COLUMNWIDTH\PY{n+nb}{\PYZus{}}PT=\PY{k}{\PYZbs{}the}\PY{k}{\PYZbs{}columnwidth}\PY{n+nb}{\PYZcb{}}
\PY{k}{\PYZbs{}typeout}\PY{n+nb}{\PYZob{}}PAPERWIDTH\PY{n+nb}{\PYZus{}}PT=\PY{k}{\PYZbs{}the}\PY{k}{\PYZbs{}paperwidth}\PY{n+nb}{\PYZcb{}}
\PY{k}{\PYZbs{}typeout}\PY{n+nb}{\PYZob{}}PAPERHEIGHT\PY{n+nb}{\PYZus{}}PT=\PY{k}{\PYZbs{}the}\PY{k}{\PYZbs{}paperheight}\PY{n+nb}{\PYZcb{}}
\PY{k}{\PYZbs{}thispagestyle}\PY{n+nb}{\PYZob{}}empty\PY{n+nb}{\PYZcb{}}

\PY{k}{\PYZbs{}sbox}\PY{n+nb}{\PYZob{}}\PY{k}{\PYZbs{}tblbox}\PY{n+nb}{\PYZcb{}}\PY{n+nb}{\PYZob{}}\PY{c}{\PYZpc{}}
  \PY{k}{\PYZbs{}begin}\PY{n+nb}{\PYZob{}}tabular\PY{n+nb}{\PYZcb{}}\PY{n+nb}{\PYZob{}}l l r\PY{n+nb}{\PYZcb{}}
    \PY{k}{\PYZbs{}toprule}
    Item \PY{n+nb}{\PYZam{}} Description \PY{n+nb}{\PYZam{}} Price\PY{k}{\PYZbs{}\PYZbs{}}
    \PY{k}{\PYZbs{}midrule}
    Pink Rabbit \PY{n+nb}{\PYZam{}} A small pink puppet made of soft fur. \PY{n+nb}{\PYZam{}} \PY{k}{\PYZbs{}\PYZdl{}}10\PY{k}{\PYZbs{}\PYZbs{}}
    \PY{k}{\PYZbs{}bottomrule}
  \PY{k}{\PYZbs{}end}\PY{n+nb}{\PYZob{}}tabular\PY{n+nb}{\PYZcb{}}\PY{c}{\PYZpc{}}
\PY{n+nb}{\PYZcb{}}

\PY{k}{\PYZbs{}typeout}\PY{n+nb}{\PYZob{}}TABLE\PY{n+nb}{\PYZus{}}WD\PY{n+nb}{\PYZus{}}PT=\PY{k}{\PYZbs{}the}\PY{k}{\PYZbs{}wd}\PY{k}{\PYZbs{}tblbox}\PY{n+nb}{\PYZcb{}}
\PY{k}{\PYZbs{}typeout}\PY{n+nb}{\PYZob{}}TABLE\PY{n+nb}{\PYZus{}}HT\PY{n+nb}{\PYZus{}}PT=\PY{k}{\PYZbs{}the}\PY{k}{\PYZbs{}dimexpr}\PY{k}{\PYZbs{}ht}\PY{k}{\PYZbs{}tblbox}+\PY{k}{\PYZbs{}dp}\PY{k}{\PYZbs{}tblbox}\PY{k}{\PYZbs{}relax}\PY{n+nb}{\PYZcb{}}

\PY{k}{\PYZbs{}end}\PY{n+nb}{\PYZob{}}document\PY{n+nb}{\PYZcb{}}
\end{Verbatim}
\end{tscode}
\end{document}
